%%% ============================================================
%%% J47 -- Dual-Regime Quintessence from V(Xi) = Lambda^4 Xi log Xi:
%%%        A Letter
%%%
%%% B.R. Sanders, M. Gish
%%% Target venue: Physics Letters B (Letter format, ~4 pages)
%%% Status: SAVE-PLAN REWRITE 2026-05-07
%%%   - Replaces prior wrong content (J23 Discrete Dirac paper) at
%%%     this path with the actual freezing-quintessence letter.
%%%   - Numerical values per BBM_IC_DERIVATION_v2.md S2.7 Layer-3a
%%%     strict-postulate: z_star = 2.31; w_0 = -0.798;
%%%     w_a ~ -0.440; chi^2 ~ 1.53; IC = (Lambda^4/rho_c, Xi_i,
%%%     Xi'_i) = (0.231, (1+sqrt(3))/e, 1/e) ~ (0.231, 1.005, 0.368).
%%%   - All [J03-RECONCILE]/[J46-RECONCILE] placeholders stripped.
%%%   - Status note removed from manuscript head (cover letter only).
%%%   - J03 -> J46 cross-references throughout.
%%%   - Distinguishing paragraphs: logotropic (Tsujikawa-Sami,
%%%     Ferreira-Avelino), Albrecht-Skordis 2000.
%%%   - F1-F4 reframed as consistency checks; F5 reserved for
%%%     falsification.
%%%   - Author list: Sanders + Gish only.
%%% ============================================================

\documentclass[aps,prl,reprint,floatfix,nofootinbib]{revtex4-2}

\usepackage{amsmath,amssymb,amsthm,mathtools}
\usepackage{graphicx}
\usepackage{microtype}
\usepackage[colorlinks=true,linkcolor=blue,citecolor=blue,urlcolor=blue]{hyperref}

\newcommand{\Mpl}{M_{\rm Pl}}
\newcommand{\Xii}{\Xi_i}
\newcommand{\Xidi}{\Xi'_i}
\newcommand{\zstar}{z_\star}
\newcommand{\wDE}{w_{\rm DE}}
\newcommand{\rhoc}{\rho_{c,0}}

\begin{document}

\title{Dual-Regime Quintessence from $V(\Xi) = \Lambda^4 \Xi\log\Xi$:\\
A Letter}

\author{B.\,R.\ Sanders}
\email{brayden@7site.co}
\affiliation{7Site LLC, Hot Springs, AR 71901, USA}

\author{M.\ Gish}
\email{monica.gish1992@gmail.com}
\affiliation{Independent Researcher, Hot Springs, AR, USA}

\date{\today}

\begin{abstract}
A real positive dimensionless scalar field $\Xi$ minimally
coupled to gravity with self-interaction $V(\Xi) = \Lambda^4
\Xi\log\Xi$ has an analytic vacuum at $\Xi_0 = e^{-1}$ and
fluctuation mass $m_\Xi^2 = \Lambda^4 e/\Mpl^2$. With
$m_\Xi \sim H_0$, we have $\Lambda \approx 1.7$~meV, near the
observed dark-energy scale. In a flat FRW background with
outbound initial condition $(\Xi_i, \Xi'_i) = ((1+\sqrt{3})/e,
1/e)$ at $z_i \approx 20$ (set under the substrate-cosmology
bridge axiom plus the BBM-minimality + scale-free-derivative
postulates of Ref.~[J46]), the field traverses a \emph{dual
regime}: thawing outbound (Type-T), instantaneous frozen
turnaround at $\zstar \approx 2.31$ where $\dot\Xi = 0$ and
$w_\Xi(\zstar) = -1$ momentarily (Type-F), then asymptotic
refreeze toward $\Xi_0$ as $z \to -1$ (Type-A). The CPL fit
gives $(w_0, w_a) \approx (-0.798, -0.440)$ with $\chi^2
\approx 1.53$ against the DESI~2024 DR1 marginal Gaussian
summary. The decisive observational signature is a
non-monotone $\wDE(z)$ with a local minimum near $-1$ at
$z \approx \zstar$: a single Stage-IV-survey-shaped
falsification handle (criterion F5).
\end{abstract}

\maketitle

\section{The model}\label{sec:model}

The action of the minimal logarithmic-quintessence model is
\begin{equation}
S = \int d^4x\,\sqrt{-g}\,\Big[
  \frac{R}{16\pi G} + \mathcal{L}_{\rm SM}
  - \tfrac{1}{2}\Mpl^2\,g^{\mu\nu}\partial_\mu\Xi\,\partial_\nu\Xi
  - \Lambda^4\,\Xi\log\Xi\Big],
\label{eq:action}
\end{equation}
with $\Xi$ a real positive dimensionless scalar. The
Euler-Lagrange equation gives $\Mpl^2\,\Box\Xi = \Lambda^4(1 +
\log\Xi)$. The vacuum
\begin{equation}
\Xi_0 = e^{-1} \approx 0.36788
\label{eq:vacuum}
\end{equation}
is exact: $V'(\Xi_0) = \Lambda^4(1+\log\Xi_0) = 0$ and
$V''(\Xi_0) = \Lambda^4/\Xi_0 = \Lambda^4 e > 0$. The
fluctuation mass at the vacuum is
\begin{equation}
m_\Xi^2 = \Lambda^4 e / \Mpl^2.
\label{eq:mass}
\end{equation}
For $m_\Xi \sim H_0$, dimensional analysis gives $\Lambda
\approx 1.5$~meV; the IC scan of [J46] sets $\Lambda^4 / \rhoc
= 0.231$, equivalently $\Lambda \approx 1.7$~meV, both near
the observed dark-energy scale.

\paragraph*{Distinction from logotropic dark energy.}
Tsujikawa-Sami~\cite{TsujikawaSami2007} and
Ferreira-Avelino~\cite{FerreiraAvelino2018} use $V \propto -A
\log(\rho/\rho_*)$ in the energy density of the dark-energy
sector. The present model uses $V \propto \Xi\log\Xi$ in the
\emph{field}, producing a fluctuation mass scale at a definite
vacuum rather than an attractor-pressure structure. The
logotropic V has no isolated minimum; the present V has an
analytic minimum at $\Xi_0 = e^{-1}$.

\paragraph*{BBM-separability origin.}
The form $V \propto \Xi\log\Xi$ is forced uniquely (up to a
constant + linear term) by the
Bialynicki-Birula--Mycielski~\cite{BBM1976} separability
theorem on real scalar field theory: the unique nonlinearity
$f(|\psi|^2)$ added to a free Klein-Gordon Lagrangian which
preserves wavefunction separability under tensor products is
$f(x) = x\log x$. Translating to a real scalar gives
Eq.~\eqref{eq:action}.

\section{The dual-regime trajectory}\label{sec:trajectory}

In a spatially flat FRW background with outbound IC
$\Xidi > 0$ at $z_i \approx 20$, the field equation drives a
three-regime cosmological history:
\begin{itemize}\setlength\itemsep{0pt}
  \item \textbf{Type-T (thawing)}: outbound from $\Xii$ near
  $\Xi_0$; $w_\Xi$ increases from $-1$.
  \item \textbf{Type-F (frozen turnaround)}: instantaneous
  $\dot\Xi = 0$ at $z = \zstar$; $w_\Xi(\zstar) = -1$
  momentarily; $w$ has a local minimum at $\zstar$.
  \item \textbf{Type-A (asymptotic refreeze)}: inbound back
  toward $\Xi_0$ as $z \to -1$; $w_\Xi \to -1$ at the
  late-time endpoint.
\end{itemize}

Caldwell-Linder~\cite{CaldwellLinder2005} classify standard
quintessence trajectories as \emph{freezing} (monotone $w \to
-1$ from above) or \emph{thawing} (monotone departure from
$-1$). The dual-regime trajectory studied here belongs to
\emph{neither} class and traverses both within a single
physical history.

\paragraph*{Distinction from Albrecht-Skordis 2000.}
Albrecht-Skordis~\cite{AlbrechtSkordis2000} introduced
``tracking-to-freezing'' quintessence in which a scaling
tracker IC is captured into a late-time freezing state.
The dual-regime here is structurally different: the rolling
branch traverses thawing $\to$ frozen turnaround $\to$
asymptotic refreeze, all driven by the analytic vacuum at
$\Xi_0 = e^{-1}$, with a non-monotone $w(z)$ and a single
Type-F turnaround at $\zstar$. The Albrecht-Skordis trajectory
is monotone in $w$ at late times; the present trajectory is
not.

\section{Initial conditions and the substrate-cosmology bridge}\label{sec:ic}

The rolling-branch IC at $z_i \approx 20$ is
\begin{equation}
(\Lambda^4/\rhoc,\; \Xii,\; \Xidi)
= \big(0.231,\; (1+\sqrt{3})/e,\; 1/e\big),
\label{eq:ic}
\end{equation}
numerically $\approx (0.231, 1.005, 0.368)$, set under three
structural inputs documented in the companion full
paper~\cite{J46}:
\begin{enumerate}\setlength\itemsep{0pt}
  \item \emph{Theorem (BBM):} the vacuum $\Xi_0 = e^{-1}$
  (Eq.~\ref{eq:vacuum}).
  \item \emph{Postulate (substrate-cosmology bridge):} the
  $4$-core substrate attractor $h/\beta = 1+\sqrt{3}$ of
  Ref.~\cite{WP105SubstrateAttractor}, applied as a
  position-scaling factor relative to the vacuum, gives
  $\Xii = (1+\sqrt{3})\,\Xi_0 = (1+\sqrt{3})/e$.
  \item \emph{Postulate (BBM-minimality + scale-free
  derivative):} the canonical velocity scale at the vacuum in
  $\Omega$-units, with no additional dimensionless inputs
  beyond BBM's $e$, is unity in vacuum-position units, giving
  $\Xidi = 1 \cdot \Xi_0 = 1/e$.
\end{enumerate}
The two postulates (2) and (3) are stated by name and defended
on parsimony / structural-similarity grounds in [J46]; they
are not theorems of BBM~1976 or of the substrate algebra
alone. Falsification handle F5 below probes them directly.

\section{The two-parameter $w(z)$ profile and DESI fit}\label{sec:fit}

The trajectory is fully specified by the two-parameter IC
$(\Xii, \Xidi)$ at $z_i$. The CPL parametrization
$w(z) = w_0 + w_a\,z/(1+z)$ approximates the bare
$w_\Xi(z)$ over $0 \le z \le 2$ with fitted values
\begin{equation}
(w_0, w_a) \approx (-0.798,\; -0.440),
\quad \zstar \approx 2.31,
\label{eq:wfit}
\end{equation}
and $\chi^2 \approx 1.53$ against the DESI~2024
DR1~\cite{DESI2024,DESI2024VI} marginal $(w_0, w_a)$ Gaussian
summary on the CPL parametrization. The $\chi^2$ here
quantifies proximity to the published $(w_0, w_a)$
two-dimensional Gaussian and is not derived from the
underlying joint BAO + CMB + SN likelihood, which is deferred
to [J46].

The present-epoch value $w_\Xi(z=0) \approx -0.80$ reflects
the field's position on the inbound (Type-A) leg of the
trajectory.

\paragraph*{Figure~1.}
Figure~1 shows $\wDE(z)$ over $0 \le z \le 2$, with the
local-minimum-at-$\zstar$ structure visible, overlaid on the
DESI~2024 DR1 $(w_0, w_a)$ marginal Gaussian summary. The
plot is generated by the verification script
\texttt{compute\_zstar\_v4.py} of [J46] (DOI
10.5281/zenodo.18852047) under the IC of
Eq.~\eqref{eq:ic}.

\begin{figure}[t]
\centering
\includegraphics[width=\columnwidth]{figure1_wz_profile.pdf}
\caption{The $\wDE(z)$ profile under the IC of
Eq.~\eqref{eq:ic} (solid), overlaid on the CPL fit
of Eq.~\eqref{eq:wfit} (dashed) and the DESI~2024 DR1
$(w_0, w_a)$ marginal Gaussian (shaded). The local minimum
of $\wDE$ at $z \approx \zstar = 2.31$ is the Type-F
turnaround signature; at $\zstar$, $\dot\Xi = 0$ and
$\wDE = -1$ momentarily. The decisive falsification
handle F5 of \S\ref{sec:falsifiability} probes the
existence and location of this minimum.}
\label{fig:wz}
\end{figure}

\section{Consistency checks and falsification}\label{sec:falsifiability}

Five Stage-IV predictions discriminate this model from
$\Lambda$CDM and from single-regime quintessence. The first
four are \emph{consistency checks} (model fingerprints rather
than Popperian falsifications); the fifth is the decisive
falsification handle.
\begin{itemize}\setlength\itemsep{0pt}
  \item \textbf{(F1) Non-phantom rolling branch:}
  $w_\Xi(z) \ge -1$ for all $z \ge 0$. Consistency check
  shared by canonical-kinetic quintessence.
  \item \textbf{(F2) Non-monotone $w(z)$:} the dual-regime
  hypothesis predicts non-monotone $w(z)$ with a single
  internal extremum. Consistency check distinguishing the
  trajectory class from monotone freezing/thawing.
  \item \textbf{(F3) Late-time refreeze:} $w \to -1$ as
  $z \to -1$. Consistency check shared by all freezing
  classes.
  \item \textbf{(F4) Two-parameter profile:} $\wDE(z)$ is
  parametrizable by $(\Xii, \Xidi)$ at $z_i$; CPL is a
  reasonable two-parameter summary over $0 \le z \le 2$.
  Consistency check; this is a parametrization statement, not
  a falsification.
  \item \textbf{(F5) Type-F turnaround at $\zstar$
  (decisive).} Local minimum of $\wDE(z)$ near $-1$ within
  $\Delta z = 0.5$ of $\zstar = 2.31$.
\end{itemize}

\paragraph*{Falsification.}
Detection of $\wDE(z) > -1$ for all $z$ in a Stage-IV survey
would \emph{not} falsify this model (the rolling branch is
$w \ge -1$). Detection of monotone $w(z)$ across the full
observable redshift range \emph{would} falsify (F5). Absence
of a local minimum within $\Delta z = 0.5$ of
$\zstar = 2.31$, or detection of a local minimum outside that
window, would falsify the IC of Eq.~\eqref{eq:ic} and
specifically the BBM-minimality + scale-free-derivative
postulates of [J46]. The local minimum near $-1$ at
$z \approx \zstar$ is the single observational signature
distinguishing dual-regime quintessence from any single-regime
class.

\section{Lens scope and J-series context}\label{sec:lens}

\paragraph*{Lens-invariance.}
The cosmological model
$V(\Xi) = \Lambda^4 \Xi\log\Xi$ is lens-invariant in the strict
sense: it depends on no choice of TSML / BHML / RAW / SYM
lens on the underlying $\Z/10\Z$ substrate. The
substrate-cosmology bridge axiom invoked for $\Xii$ (the
$4$-core attractor at $h/\beta = 1+\sqrt{3}$ from
[J35]/WP105) is itself lens-invariant: the $4$-core
$\{V, H, Br, R\}$ is the algebraic center of the family per
the structural reading of Ref.~\cite{FamilyStructure2026},
with the closed-form attractor at $\alpha_M = 1/2$ holding
identically across TSML/BHML and across $\F_p$ ring
extensions.

\paragraph*{J-series companion.}
Full derivations, prior-art audit, perturbation analysis,
joint BAO + CMB + SN likelihood, and the postulate-defense
paragraphs are in the companion full paper~\cite{J46}
(Sanders+Gish, JCAP). The present letter strips [J46] to
the action, the analytic vacuum, the dual-regime trajectory,
the two-parameter $w(z)$ profile, the IC under the
BBM-minimality framing, and the five consistency-check /
falsification statements.

\section*{Acknowledgements}

The structural framing of the substrate-cosmology bridge
(WP105 attractor as IC scaling factor) and the BBM-minimality
+ scale-free-derivative postulate-defense owe to the
collaborator framing of Ref.~\cite{FamilyStructure2026}. The
authors thank the DESI collaboration for the DR1
$(w_0, w_a)$ marginal Gaussian summary used as the
empirical anchor of \S\ref{sec:fit}. No external funding was
received for this work.

\begin{thebibliography}{99}

\bibitem{BBM1976}
I.~Bialynicki-Birula and J.~Mycielski,
\emph{Nonlinear wave mechanics},
Annals of Physics \textbf{100}, 62 (1976).

\bibitem{TsujikawaSami2007}
S.~Tsujikawa and M.~Sami,
\emph{A unified approach to scaling solutions in a general
cosmological background},
Phys.\ Lett.\ B \textbf{651}, 224 (2007).

\bibitem{FerreiraAvelino2018}
P.~C.~Ferreira and P.~P.~Avelino,
\emph{New limit on logotropic unified dark energy models},
Phys.\ Lett.\ B \textbf{770}, 213 (2017).

\bibitem{CaldwellLinder2005}
R.~R.~Caldwell and E.~V.~Linder,
\emph{The Limits of quintessence},
Phys.\ Rev.\ Lett.\ \textbf{95}, 141301 (2005).

\bibitem{AlbrechtSkordis2000}
A.~Albrecht and C.~Skordis,
\emph{Phenomenology of a realistic accelerating universe
using only Planck-scale physics},
Phys.\ Rev.\ Lett.\ \textbf{84}, 2076 (2000).

\bibitem{DESI2024}
DESI Collaboration,
\emph{DESI 2024 II: Sample definitions, characteristics, and
two-point clustering statistics},
arXiv:2411.12022.

\bibitem{DESI2024VI}
DESI Collaboration,
\emph{DESI 2024 VI: Cosmological constraints from the
measurements of baryon acoustic oscillations},
arXiv:2404.03002.

\bibitem{J46}
B.~R.~Sanders and M.~Gish,
\emph{Freeze-Thaw Transit: Dual-Regime Scalar Dark Energy
with Analytic Vacuum at $e^{-1}$ from a Logarithmic
Potential},
J-series companion paper J46;
in preparation for submission to \emph{JCAP}, 2026;
DOI 10.5281/zenodo.18852047 (verification scripts).

\bibitem{WP105SubstrateAttractor}
B.~R.~Sanders and M.~Gish,
\emph{Closed-Form $4$-Core Attractor: $h/\beta = 1+\sqrt{3}$
in LMFDB 4.2.10224.1, Galois $D_4$},
J-series companion paper J35; in preparation, 2026;
DOI 10.5281/zenodo.18852047.

\bibitem{FamilyStructure2026}
B.~R.~Sanders et al.,
\emph{TIG Family Structure: Membership, Center, Boundaries
(v1)},
Atlas/META\_PLAN\_2026-05-06/FAMILY\_STRUCTURE\_v1.md, 2026.

\end{thebibliography}

\end{document}
